Allow us to define the mathematical objects we utilize throughout the manuscript.
For multifold sums of powers, we use the notation provided by Donald Knuth in~\cite{knuth1993johann},
because it is simple and beautiful.
\begin{definition} [Multifold sums of powers]
  \label{def:multifold-sums-of-powers-recurrence}
  For non-negative integers $r,n,m$
  \begin{align*}
    \KnuthRFoldSum{0}{n}{m}   &= n^m, \\
    \KnuthRFoldSum{1}{n}{m}   &= \KnuthRFoldSum{0}{1}{m}
    + \KnuthRFoldSum{0}{2}{m} + \cdots + \KnuthRFoldSum{0}{n}{m}, \\
    \KnuthRFoldSum{r+1}{n}{m} &= \KnuthRFoldSum{r}{1}{m}
    + \KnuthRFoldSum{r}{2}{m} + \cdots + \KnuthRFoldSum{r}{n}{m}.
  \end{align*}
\end{definition}
We utilize the following notation for falling factorials
\begin{definition} [Falling factorials] For integers $x,n$
  \begin{align*}
    \fallingFactorial{x}{n} =
    \begin{cases}
      1, & \mathrm{if \; } k = 0, \\
      x(x-1)(x-2)\cdots(x-n+1) = \prod_{k=0}^{n-1} (x-k), & \mathrm{if \; } k > 0, \\
      0, & \mathrm{else}.
    \end{cases}
  \end{align*}
\end{definition}
We utilize the following notation for central factorials
\begin{definition}[Central factorials]
  \label{def:central-factorials}
  For integers $n,k$
  \begin{align*}
    \centralFactorial{n}{k} =
    \begin{cases}
      0, \quad & \mathrm{if \; } k < 0, \\
      1, \quad & \mathrm{if \; } k = 0, \\
      n \left( n + \frac{k}{2} -1 \right)\left( n + \frac{k}{2} -2 \right)
      \cdots \left( n - \frac{k}{2} +1 \right)
      = n \fallingFactorial{n+\frac{k}{2}-1}{k-1}, \quad & \mathrm{if \; } k>0,
    \end{cases}
  \end{align*}
  where
  $\fallingFactorial{n+\frac{k}{2}-1}{k-1}
  = \prod_{j=0}^{k-2} \left( n+\frac{k}{2}-1 -j \right)$
  are falling factorials.
\end{definition}
We encourage the audience to read
J. F. Steffensen's work on the definition of
generalized factorials~\cite{steffensen1933definition}.
